\chapter{Exponencial y Logaritmo}\label{ch:g12:exp}

El \omterm{def:g11:func:function}{función} exponencial es el unico \omterm{def:g11:func:function}{función} igual a su propio derivada
y tomando el valor $1$ en $0$. Convierte sumas en productos; es inverso,
el logaritmo natural, convierte productos en sumas. Juntos describen
Todo fenómeno cuya tasa de cambio es proporcional a su tamaño:
desintegración radiactiva, crecimiento demográfico, interés compuesto.

\section{La función exponencial}

\begin{theorem}[Existencia y unicidad]\label{thm:g12:exp:existence}
Existe un único \omterm{def:g12:deriv:derivative}{diferenciable} \omterm{def:g11:func:function}{función} $ f \colon \R \to \R $ tal que
\[
f' = f \qquad\text{and}\qquad f(0) = 1 .
\]
Se llama el \emph{funcion exponencial}\index{funcion exponencial} y escrito
$\exp $, o $ x \mapsto \eu^x $.
\end{theorem}

\begin{proof}[Prueba de unicidad]
Primero, tal \omterm{def:g11:func:function}{función} nunca desaparece. De hecho, deja
$ g(x) = f(x)f(-x)$; entonces
\[
g'(x) = f'(x)f(-x) - f(x)f'(-x) = f(x)f(-x) - f(x)f(-x) = 0,
\]
entonces $ g $ es constante igual a $ g(0) = 1$: por cada $ x $,
$ f(x)f(-x) = 1$, y en particular $ f(x) \neq 0$.

Ahora deja $ f_1, f_2$ ser dos soluciones y $ h = \dfrac{f_1}{f_2}$ (legítimo
desde $ f_2$ nunca desaparece). Entonces
\[
h' = \frac{f_1' f_2 - f_1 f_2'}{f_2^2} = \frac{f_1 f_2 - f_1 f_2}{f_2^2} = 0,
\]
entonces $ h $ es constante igual a $ h(0) = 1$, \emph{es decir.} $ f_1 = f_2$.

\emph{Existencia} está admitido en este nivel (se puede obtener a través del
\omterm{def:g12:seq:sequence}{sucesión} de \cref{exo:g12:seq:10}, o como la inversa de la logarítmo
construido por integración en \cref{ch:g12:integ}).
\end{proof}

\begin{proposition}[Ecuación funcional]\label{prop:g12:exp:funceq}
A pesar de $ x, y \in \R $ y $ n \in \Z $:
\[
\eu^{x+y} = \eu^x\,\eu^y, \qquad
\eu^{-x} = \frac{1}{\eu^x}, \qquad
\eu^{x-y} = \frac{\eu^x}{\eu^y}, \qquad
\bigl(\eu^{x}\bigr)^n = \eu^{nx}.
\]
Además $\eu^x > 0$ para todo $ x $.
\end{proposition}

\begin{proof}
Arreglar $ y $ y considerar $\varphi(x) = \dfrac{\exp(x+y)}{\exp(x)\exp(y)}$. su
numerador y denominador, como funciones de $ x $, son ambas soluciones de
$ f' = f $ hasta la constante $\exp(y)$; diferenciando $\varphi $ directamente
(regla del cociente) da $\varphi' = 0$, entonces
$\varphi \equiv \varphi(0) = 1$, demostrando la primera identidad. Tomando
$ y = -x $ da el segundo (con el valor $\eu^0 = 1$), y el tercero
sigue. El cuarto es una inducción del primero para $ n \geq 0$, extendido
a $ n < 0$ por el segundo.

Positividad: $\eu^x = \left(\eu^{x/2}\right)^2 \geq 0$ y $\eu^x \neq 0$
(mostrado en \cref{thm:g12:exp:existence}), entonces $\eu^x > 0$.
\end{proof}

\begin{proposition}[Variaciones y límites]\label{prop:g12:exp:variations}
La exponencial es estrictamente \omterm{def:g12:seq:monotonic}{creciente}, convexa, y
\[
\lim_{x\to-\infty} \eu^x = 0, \qquad
\lim_{x\to+\infty} \eu^x = +\infty .
\]
\end{proposition}

\begin{proof}
$(\exp)' = \exp > 0$ da aumento estricto; $(\exp)'' = \exp > 0$ da
convexidad. Por convexidad, $\eu^x \geq 1 + x $ (\omterm{def:g12:deriv:tangent}{tangente} en $0$, ver
\cref{ex:g12:deriv:ineq}), entonces $\eu^x \to +\infty $ como $ x \to +\infty $ por
comparación. Entonces $\eu^{x} = 1/\eu^{-x} \to 0$ como $ x \to -\infty $.
\end{proof}

\begin{omfigure}
\begin{tikzpicture}
\begin{axis}[omaxis,
  width=8.8cm, height=6cm,
  xmin=-3.3, xmax=2.4, ymin=-2.4, ymax=7.8,
  xtick={-2,-1,1,2}, ytick={1}]
  \addplot[omProp, thick, domain=-3.2:2.3] {1 + x};
  \addplot[omDef, thick, domain=-3.2:2.0] {exp(x)};
  \fill (axis cs:0,1) circle (1.5pt);
  \node[omDef, font=\small, anchor=east] at (axis cs:1.35,5.2) {$\eu^x $};
  \node[omProp, font=\small, anchor=north west] at (axis cs:1.05,1.6) {$ y = 1+x $};
\end{axis}
\end{tikzpicture}

{\small El exponencial: estrictamente \omterm{def:g12:seq:monotonic}{creciente}, convexa, con
$\lim_{-\infty} \eu^x = 0$. Se encuentra encima de su \omterm{def:g12:deriv:tangent}{tangente} en $0$:
$\eu^x \geq 1 + x $.}
\end{omfigure}

\begin{theorem}[Comparación de crecimiento]\label{thm:g12:exp:growth}
\index{comparación de crecimiento}
por cada entero $ n \geq 1$,
\[
\lim_{x\to+\infty} \frac{\eu^x}{x^n} = +\infty,
\qquad
\lim_{x\to-\infty} x^n \eu^x = 0 .
\]
En palabras: la exponencial supera a todas las potencias de $ x $.
\end{theorem}

\begin{proof}
Para $ n = 1$: aplicando $\eu^t \geq 1 + t \geq t $ en $ t = x/2$,
\[
\frac{\eu^x}{x} = \frac{\left(\eu^{x/2}\right)^2}{x}
\geq \frac{(x/2)^2}{x} = \frac{x}{4} \xrightarrow[x\to+\infty]{} +\infty .
\]
para generales $ n $, escribir
\[
\frac{\eu^x}{x^n}
= \frac{1}{n^n}\left(\frac{\eu^{x/n}}{x/n}\right)^{n}
\xrightarrow[x\to+\infty]{} +\infty
\]
por el caso $ n=1$ y composición. El límite en $-\infty $ sigue por el
sustitución $ x \mapsto -x $:
$\abs{x^n \eu^x} = \frac{\abs{x}^n}{\eu^{\abs x}} \to 0$.
\end{proof}

\begin{omfigure}
\begin{tikzpicture}
\begin{axis}[omaxis,
  width=8.8cm, height=6cm,
  xmin=-0.4, xmax=5.6, ymin=-12, ymax=155,
  xtick={1,2,3,4,5}, ytick={50,100,150}]
  \addplot[omThm, thick, domain=0:5.4] {x^2};
  \addplot[omProp, thick, domain=0:5.3] {x^3};
  \addplot[omDef, thick, domain=-0.4:5.0] {exp(x)};
  \fill (axis cs:4.536,93.4) circle (1.5pt);
  \draw[black, dashed, thin] (axis cs:4.536,0) -- (axis cs:4.536,93.4);
\end{axis}
\end{tikzpicture}

{\small El exponencial (azul) finalmente vence a todas las potencias:
adelanta $ x^2$ (rojo) desde el principio y $ x^3$ (naranja) en $ x \approx 4.54$.}
\end{omfigure}

\section{El logaritmo natural}

\begin{definition}[logaritmo natural]\label{def:g12:exp:ln}
La exponencial es continua y estrictamente \omterm{def:g12:seq:monotonic}{creciente} de $\R $ sobre
$\intoo{0}{+\infty}$; por el teorema de biyección
(\cref{thm:g12:limcont:bijection}), por cada $ y > 0$ el \omterm{def:g10:algebra:equation}{ecuación}
$\eu^x = y $ tiene una solución única. Esta solución es la \emph{naturales
logarítmo}\index{naturales
logarítmo} de $ y $, escrito $\ln y $. De este modo
\[
\text{for } x \in \R,\ y > 0: \qquad y = \eu^x \iff x = \ln y .
\]
En particular $\ln 1 = 0$, $\ln \eu = 1$, y
$\eu^{\ln y} = y $, $\ln(\eu^x) = x $.
\end{definition}

\begin{omfigure}
\begin{tikzpicture}
\begin{axis}[omaxis,
  width=8.6cm, axis equal image,
  xmin=-4.3, xmax=6.3, ymin=-4.3, ymax=6.3,
  xtick={1}, ytick={1}]
  \addplot[black, dashed, domain=-3.9:5.9] {x};
  \addplot[omDef, thick, domain=-4.1:1.79] {exp(x)};
  \addplot[omThm, thick, domain=0.0167:5.9, samples=120] {ln(x)};
  \node[omDef, font=\small, anchor=east] at (axis cs:0.95,2.6) {$\eu^x $};
  \node[omThm, font=\small, anchor=north] at (axis cs:4.3,1.25) {$\ln x $};
  \node[font=\small, anchor=west, rotate=45] at (axis cs:4.6,4.85) {$ y=x $};
\end{axis}
\end{tikzpicture}

{\small las curvas de $\exp $ y $\ln $ son espejo imágenes el uno del otro
en la linea $ y = x $: el punto $(x, \eu^x)$ refleja a $(\eu^x, x)$.}
\end{omfigure}

\begin{proposition}[Propiedades algebraicas]\label{prop:g12:exp:lnalg}
A pesar de $ a, b > 0$ y $ n \in \Z $:
\[
\ln(ab) = \ln a + \ln b, \quad
\ln\frac{1}{a} = -\ln a, \quad
\ln\frac{a}{b} = \ln a - \ln b, \quad
\ln(a^n) = n \ln a, \quad
\ln\sqrt{a} = \tfrac12 \ln a .
\]
\end{proposition}

\begin{proof}
$\eu^{\ln a + \ln b} = \eu^{\ln a}\eu^{\ln b} = ab $, y tomando $\ln $ de ambos
lados da la primera identidad. Los demás siguen por el mismo mecanismo desde
las identidades correspondientes de \cref{prop:g12:exp:funceq}.
\end{proof}

\begin{proposition}[Propiedades analíticas]\label{prop:g12:exp:lnanalytic}
La \omterm{def:g11:func:function}{función} $\ln $ es \omterm{def:g12:deriv:derivative}{diferenciable} en $\intoo{0}{+\infty}$ con
\[
(\ln x)' = \frac{1}{x},
\]
estrictamente \omterm{def:g12:seq:monotonic}{creciente}, cóncava, y
$\lim\limits_{x\to0^+} \ln x = -\infty $,
$\lim\limits_{x\to+\infty} \ln x = +\infty $. Es más, por cada entero
$ n \geq 1$,
\[
\lim_{x\to+\infty} \frac{\ln x}{x^{1/n}} = 0, \qquad
\lim_{x\to0^+} x \ln x = 0 .
\]
\end{proposition}

\begin{proof}
Diferenciabilidad de la inversa \omterm{def:g11:func:function}{función} es admitido en este nivel; otorgando
diferenciar la identidad $\eu^{\ln x} = x $ por la regla de la cadena:
$(\ln x)'\,\eu^{\ln x} = 1$, entonces $(\ln x)' = \frac{1}{x} > 0$, de donde estricto
aumentar, y $(\ln x)'' = -\frac1{x^2} < 0$, de donde la concavidad. Los límites en
$0^+$ y $+\infty $ reflejar los de $\exp $ a través de la biyección. Para el
\omterm{thm:g12:exp:growth}{comparación de crecimiento}, sustituto $ x = \eu^{t}$:
$\frac{\ln x}{x} = \frac{t}{\eu^t} \to 0$ como $ t \to +\infty $ por
\cref{thm:g12:exp:growth}, y de manera similar para los otros límites con
$ x = \eu^{-t}$, donación $ x\ln x = -t\eu^{-t} \to 0$.
\end{proof}

\begin{method}[Resolver ecuaciones con $\exp $ y $\ln $]\label{met:g12:exp:equations}
Ambos funciones son estrictamente \omterm{def:g12:seq:monotonic}{creciente}, para que puedan aplicarse (o eliminarse)
de) ambos lados de un \omterm{def:g10:algebra:equation}{ecuación} o desigualdad sin cambiar su dirección:
\[
\eu^{u} = \eu^{v} \iff u = v, \qquad
\eu^{u} \leq \eu^{v} \iff u \leq v,
\]
y de la misma manera para $\ln $ en cantidades positivas. \emph{Siempre revisa \omterm{def:g11:func:function}{dominios}
primero} ($\ln $ requiere argumentos positivos). Para ecuaciones como
$ a^x = b $ con $ a>0$, $ a \neq 1$, reescribir $ a^x = \eu^{x\ln a}$ y resolver
$ x = \frac{\ln b}{\ln a}$.
\end{method}

\begin{example}\label{ex:g12:exp:halflife}
Un radiactivo \omterm{def:g10:proba:sample}{muestra} decae siguiendo $ N(t) = N_0\,\eu^{-\lambda t}$. Es
\emph{vida \omterm{def:g11:stat:mean}{media}} $ T $ satisface $ N(T) = N_0/2$, \emph{es decir.}
$\eu^{-\lambda T} = \frac12$, entonces $ T = \frac{\ln 2}{\lambda}$: la vida \omterm{def:g11:stat:mean}{media}
No depende de la cantidad inicial.
\end{example}

\section{Ejercicios}

\begin{exercise}[$\star $]\label{exo:g12:exp:1}
Simplificar
$\dfrac{\eu^{3x}\,\eu^{-x+1}}{\eu^{x}}$ y
$\ln\!\left(\dfrac{\eu^2\sqrt{\eu}}{\eu^{-3}}\right)$.
\end{exercise}

\begin{exercise}[$\star $]\label{exo:g12:exp:2}
resolver en $\R $:
\[
\text{(a) } \eu^{2x} - 3\eu^{x} + 2 = 0;
\qquad
\text{(b) } \ln(x - 1) + \ln(x + 2) = \ln 4 .
\]
\end{exercise}

\begin{exercise}[$\star $]\label{exo:g12:exp:3}
Calcular los límites:
\[
\lim_{x\to+\infty} \frac{\eu^x - x^2}{\eu^x + 1}, \qquad
\lim_{x\to+\infty} \frac{\ln x}{\sqrt x}, \qquad
\lim_{x\to0^+} x^2 \ln x, \qquad
\lim_{x\to+\infty} \bigl(x - \ln x\bigr).
\]
\end{exercise}

\begin{exercise}[$\star\star $]\label{exo:g12:exp:4}
Estudia la \omterm{def:g11:func:function}{función} $ f(x) = x\,\eu^{-x}$ en $\R $: variaciones, límites,
extremo; demuestre que su curva tiene una punto de inflexión y darle su
\omterm{def:g10:coordgeom:system}{coordenadas}.
\end{exercise}

\begin{exercise}[$\star\star $]\label{exo:g12:exp:5}
\omterm{def:g10:proba:sample}{Muestra} eso para todos $ x > -1$, $\ln(1 + x) \leq x $, con igualdad sólo en
$ x = 0$. Deduce eso para todos $ n \geq 1$,
\[
\left(1 + \frac{1}{n}\right)^n \leq \eu \leq \left(1 - \frac{1}{n+1}\right)^{-(n+1)} .
\]
\end{exercise}

\begin{exercise}[$\star\star $]\label{exo:g12:exp:6}
una capital $ C_0$ se invierte a una tasa anual de $3\%$, interés compuesto
cada año.
\begin{enumerate}
  \item expresar la capital $ C_n $ después $ n $ años.
  \item ¿Después de cuántos años se duplica el capital? da la respuesta exacta
        usando $\ln $, luego un valor numérico.
  \item Comparar con el \omterm{def:g10:numbers:approx}{aproximación} ``$70$ dividido por la tasa en
        porcentaje'' utilizado por los banqueros.
\end{enumerate}
\end{exercise}

\begin{exercise}[$\star\star $]\label{exo:g12:exp:7}
Resuelve la desigualdad $\eu^{2x} - \eu^{x+1} > 0$, entonces la desigualdad
$\ln(x^2 - 1) \leq \ln(x + 5)$.
\end{exercise}

\begin{exercise}[$\star\star\star $]\label{exo:g12:exp:8}
Sea $ f(x) = \dfrac{\eu^x}{x}$ para $ x > 0$.
\begin{enumerate}
  \item Estudiar las variaciones de $ f $ en $\intoo{0}{+\infty}$ y darle su
        \omterm{def:g10:functions:extrema}{mínimo}.
  \item ¿Para qué valores de $ k $ hace el \omterm{def:g10:algebra:equation}{ecuación} $\eu^x = kx $ tener $0$, $1$
        o $2$ soluciones en $\intoo{0}{+\infty}$?
\end{enumerate}
\end{exercise}

\begin{exercise}[$\star\star\star $]\label{exo:g12:exp:9}
Para $ n \geq 1$, dejar $ u_n = \left(1 + \frac1n\right)^n $.
\begin{enumerate}
  \item Usando \cref{exo:g12:exp:5}, \omterm{def:g10:proba:sample}{muestra} que $(u_n)$ es \omterm{def:g12:seq:bounded}{delimitado arriba}
        por ~$\eu $.
  \item mostrar que $\ln u_n = n \ln\left(1 + \frac1n\right) \to 1$, y
        deducir que $ u_n \to \eu $.
        (Pista: reconoce un cociente de \omterm{def:g11:prob:variance}{diferencias} de $\ln $ en $1$.)
\end{enumerate}
\end{exercise}

\section{Problema: el logaritmo domina al mundo}

\begin{problem}[{Problema del fin de semana --- tiempos de duplicación y la regla
de 72, las escalas de terremotos, ácidos y pianos, y la
constante $\eu $ dejando huellas dactilares por todas partes}]
\label{pb:g12:exp:1}
Todo lo que crece por un fijo \emph{porcentaje} crece
exponencialmente --- ahorros, bacterias, epidemias --- y lo que sea
abarca demasiadas potencias de diez para comprenderlas: energías sísmicas,
acidez, intensidades sonoras --- es domada por un logarítmo. esto
El problema calcula los tiempos de duplicación y desenmascara la regla de los banqueros.
72, lee las escalas logarítmicas del mundo y recopila las
huellas dactilares que el número $\eu $ deja en cada escena del crimen
(\cref{prop:g12:exp:funceq}, \cref{thm:g12:exp:growth},
\cref{met:g12:exp:equations}).

\medskip
\noindent\textbf{Part I --- Fluency.}
\begin{enumerate}
  \item Resolver: $\eu^{2x} = 5$; \quad $\ln(3x - 1) = 2$; \quad
        $\eu^{2x} - 3\eu^x + 2 = 0$ (una cuadrática disfrazada).
  \item Simplificar: $\dfrac{\ln 8}{\ln 2}$; \quad
        $\ln\!\left(\eu^3 \sqrt{\eu}\right)$; \quad
        $\eu^{\ln 5 - \ln 2}$.
  \item Demuestre la desigualdad madre: $\eu^x \geq 1 + x $ para todos
        real $ x $ (estudiar $ f(x) = \eu^x - 1 - x $), y deducir
        $\ln(1 + u) \leq u $ para $ u > -1$.
  \item Calcular
        $\lim_{x \to +\infty} x^2 \eu^{-x}$,
        $\lim_{x \to +\infty} \frac{\ln x}{x}$
        (\cref{thm:g12:exp:growth}), y
        $\lim_{x \to 0} \frac{\eu^x - 1}{x}$ (reconocer un
        derivada).
  \item Estudiar $ f(x) = x \ln x $ en $\intoo{0}{+\infty}$:
        variaciones, \omterm{def:g10:functions:extrema}{mínimo}, y el límite en $0^+$ (aceptado:
        $ x \ln x \to 0$). este pequeño \omterm{def:g11:func:function}{función} medidas
        información y entropía en los volúmenes universitarios.
\end{enumerate}

\medskip
\noindent\textbf{Part II --- Doubling times and the rule of
72.}
\begin{enumerate}[resume]
  \item Los ahorros crecen en $3\,\%$ por año. Resolver
        $1.03^n = 2$: ¿cuánto tiempo tardará en duplicarse el capital?
  \item Los banqueros estiman que el tiempo de duplicación es
        $\frac{72}{\text{rate in }\%}$. Pruebe la regla en
        $3\,\%$, $6\,\%$ y $9\,\%$ contra lo exacto
        $\frac{\ln 2}{\ln(1 + r)}$, luego explícalo: para pequeños
        $ r $, $\ln(1 + r) \approx r $ (la desigualdad de la pregunta 3 es
        la mitad de la historia), por lo que la constante exacta es
        $100 \ln 2 \approx 69.3$ --- ¿Por qué los banqueros prefieren
        $72$?
  \item Una bacteria divide cada $20$ minutos:
        $ p(t) = 2^{t/20}$ ($ t $ en minutos). Reescribirlo como
        $\eu^{\lambda t}$, luego calcular $ p $ después $24$ horas.
        La respuesta (más de $10^{21}$) demuestra ¿qué pasa con el
        modelo --- ¿y qué detiene a las verdaderas colonias?
  \item La cafeína sale del cuerpo con una vida \omterm{def:g11:stat:mean}{media} de aproximadamente $5$
        horas. Después de un $100$ mg de cafe a las 15:00, cuanto
        queda a las 23:00? ¿A qué hora baja?
        $10$ mg? (El insomnio tiene un logarítmo.)
  \item un euro en $100\,\%$ interés anual: calcular el
        Capital de fin de año anual, mensual y diario.
        composición, y dar el techo que
        \cref{exo:g12:exp:9} resultó irrompible. cual famoso
        ¿Constante es el límite de la pura codicia?
  \item ¿Por qué los científicos traman $\ln(\text{cases})$ contra
        ¿Cuál es el momento durante la fase inicial de una epidemia? ¿Qué significa un
        línea recta en esa trama revela, y ¿qué significa su
        \omterm{def:g10:reffunc:affine}{pendiente} ¿medida?
\end{enumerate}

\medskip
\noindent\textbf{Part III --- The world's logarithmic scales.}
(Escribir $\log_{10} x = \frac{\ln x}{\ln 10}$.)
\begin{enumerate}[resume]
  \item Nivel de sonido en decibeles: $ L = 10 \log_{10}(I/I_0)$. un
        medidas de conversación $60$ dB, un concierto de rock $120$ dB:
        ¿En qué factor difieren las intensidades del sonido?
  \item La magnitud de los terremotos aumenta en $1$ cuando el sísmico
        la amplitud se multiplica por $10$, y el liberado
        escalas de energía como amplitud $^{3/2}$. Comparar
        magnitud-$5$ y magnitud-$7$ terremotos: amplitud
        relación, luego relación de energía.
  \item Química: $\text{pH} = -\log_{10}[\mathrm{H^+}]$.
        El jugo de limón tiene pH. $2$, pH de la leche $6.5$: cual es el
        relación de sus concentraciones de ácido?
  \item Música: cada octava duplica la \omterm{def:g10:stats:series}{frecuencia}. Desde el
        La más baja del piano ($27.5$ Hz) a su C más alto
        ($4\,186$ Hz), ¿cuántas octavas abarca el teclado?
        ¿Y por qué nuestros sentidos (oído, vista, sensación de temblor)
        --- ¿prefieres escalas logarítmicas? (Una frase.)
  \item La regla de cálculo, la calculadora de los ingenieros para $350$
        años: dos palos graduados de manera que el \emph{longitud}
        a la marca $ x $ es proporcional a $\ln x $. Explica como
        deslizar un palo sobre el otro multiplica números,
        y nombrar la identidad de
        \cref{prop:g12:exp:lnalg} haciendo el trabajo.
\end{enumerate}

\medskip
\noindent\textbf{Part IV --- $\eu $'s fingerprints.}
\begin{enumerate}[resume]
  \item De \cref{exo:g12:exp:8}: el \omterm{def:g10:algebra:equation}{ecuación}
        $\eu^x = kx $ ($ x > 0$) tiene $0$, $1$ o $2$ soluciones
        según la posición de $ k $ relativo a un
        umbral. Vuelva a expresar el resultado --- y verifique el
        hecho geométrico detrás de esto: la línea $ y = \eu x $ es
        exactamente el \omterm{def:g12:deriv:tangent}{tangente} a la exponencial a través de la
        origen.
  \item La constante de casi accidente: con $ n $ billetes de lotería de
        ganando \omterm{def:g10:proba:distribution}{probabilidad} $\frac1n $ cada,
        $\P(\text{no win}) = \left(1 - \frac1n\right)^n $.
        Calcule su límite mediante
        $ n \ln\!\left(1 - \frac1n\right)$ (usa el estándar
        límite $\frac{\ln(1 + u)}{u} \to 1$), y reconciliarse con
        el misterioso $0.368$ de \cref{pb:g11:binom:1}.
  \item El $37\,\%$ regla: elegir lo mejor de $ n $
        candidatos entrevistados en orden aleatorio (sin vuelta atrás),
        la estrategia óptima rechaza la primera
        $\frac{n}{\eu} \approx 37\,\%$ y luego toma el primero
        candidato mejor que todos hasta ahora --- teniendo éxito con
        \omterm{def:g10:proba:distribution}{probabilidad} acerca de $\frac1\eu $. ¿Dónde están los dos?
        $\frac1\eu $ Las preguntas 18 y 19 provienen de ---
        indicar el mecanismo común (muchos independientes
        pequeña oportunidad sucesos), y calcular $\frac1\eu $ a tres
        decimales.
  \item Final --- $\eu $ El retrato: el techo de la capitalización
        (pregunta 10); la base cuya \omterm{def:g12:deriv:tangent}{tangente} en $0$ tiene \omterm{def:g10:reffunc:affine}{pendiente}
        exactamente $1$ (preguntas 3 y 4); el crecimiento no
        capturas polinómicas (pregunta 4); la constante de
        cuasi accidentes y de parada óptima (preguntas 18-19);
        y su inversa $\ln$, que convierte productos en sumas
        (pregunta 16) y décadas en pulgadas (Parte III). uno
        frase cada uno.
\end{enumerate}
\end{problem}
